% Source: 07_SGA6_EN_CURRENT_LAYERED_READER.tex:90--148.
\begin{titlepage}
\begin{center}
{\Large 博瓦马里代数几何讨论班}\par
{\large 1966--67}\par
\vspace{2em}
{\LARGE 相交理论与 Riemann--Roch 定理}\par
{\large (SGA 6)}\par
\vspace{3em}
讨论班主持人\par
P. Berthelot, A. Grothendieck, L. Illusie\par
\vspace{1em}
参与合作\par
D. Ferrand, J.-P. Jouanolou, O. Jussila,\par
S. Kleiman, M. Raynaud, J.-P. Serre
\vfill
{\large 简体中文原生重建工作稿}\par
2026 年 10 月
\end{center}
\end{titlepage}
\section*{译稿说明}
本稿是 SGA 6 的简体中文原生重建工作稿. 指定英文底稿全卷已有中文首译; 已按法文原书恢复多处被压缩的证明、脚注、公式和交换图. 2026 年 10 月 21 日, 用户要求收尾并暂停, 本项目现已暂停.

原书完整性恢复尚未完成: 第 II 讲原印页 160--191、193--208 共 48 页尚未完成法文逐页核对, 第 218 页还有一个符号待辨. 出版者投稿说明及目录广告也未重建. 逐页状态和续接点见随附的 PAUSED.md 及审计文件. 独立数学校订和全页中文视觉审查均未完成, 编译成功不表示上述项目已经通过.

各部分来自不同程度的英文转写与法文核对; 原书自身的错误保留并附译校注. 公式、交换图及正文均由原生 LaTeX 重建. 原书第十一讲未写成、未出版; 本稿保留历史说明和后续讲次编号, 不补造该讲. 原书目录页码用于对照原著.
\clearpage


\sgaStableTarget{sga6.front.preface}\section*{前言}
本版 \emph{Séminaire de Géométrie Algébrique 6} 除少量改动外, 重印了由 I.H.E.S. 以油印讲稿形式分发的第一版. 所作修正主要涉及排印错误; 只有\hyperlink{sga6.e01.expose}{第一讲}和第三讲例外, L. Illusie 对初稿中几个不正确的命题作了轻微修改. 此外还增添了一些脚注, 尤其是 A. Grothendieck 在\hyperlink{sga6.e14.expose}{第十四讲}中所加的脚注, 用来指出讨论班举行时尚未解决的问题最近取得的进展.

最后须说明, D. Ferrand 关于完美复形的行列式理论的第十一讲没有写成讲稿, 因而不会刊出; 为避免引用出错, 我们保留了后续各讲原有的编号.

\begin{flushright}1971 年 6 月\end{flushright}

\sgaStableTarget{sga6.front.introduction}\section*{引言}
本讨论班旨在为任意概形建立整体的相交理论与 Riemann--Roch 公式. 读者可在\hyperlink{sga6.e00.expose}{第零讲}中看到讨论班计划的说明. 至少对射影态射而言, 沿着著名的 1958 年 Borel--Serre 报告的思路, 为一般正则概形证明 Riemann--Roch 公式在原则上可行, 这一点当时已经清楚. 推广到一般情形则不那么明显; 这一推广所属的计划始于 1960 年, 本讨论班实现了其中的一部分. 与凝聚层的对偶理论一样, 要建立一种令人满意的理论, Verdier 的导出范畴理论是不可缺少的工具; 要理解本讨论班, 必须熟悉这一理论.

除此之外, 学习本讨论班几乎只需要对 EGA 第一、二、三章所述代数几何基础有一般了解. 我们只会偶尔引用 EGA 第四章中关于平坦态射或光滑态射的某些事实, 其中大部分也见于 SGA 1 的前几讲. 最后, 为了以应用所需的一般性发展某些结果, 我们有时使用赋环景与赋环拓扑斯的概念, 有关定义参见 SGA 4 第一至第四讲. 不熟悉景与拓扑斯语言的读者, 可以处处用通常的拓扑空间来代替这些对象, 并把拓扑斯的对象替换为空间的开子集; 不过我们仍建议读者掌握拓扑斯的语言, 因为它提供了极为方便的统一原则.

这里所述理论的基本概念是模的完美复形及其各种变体, 它们在\hyperlink{sga6.e01.expose}{第一讲}、第二讲和第三讲中展开. 显然, 这些概念在代数几何的其他领域中也很重要, 尤其是在离散系数上同调 (SGA 5) 或凝聚系数上同调的 Lefschetz--Verdier 型公式中; 它们在代数几何之外也应当发挥作用, 特别是在建立这里的 Riemann--Roch 定理的复解析变体, 或 Atiyah--Singer 指标定理的适当变体时. \hyperlink{sga6.e02.expose}{第二讲}将给出这方面的一些说明. 可惜, 目前甚至还没有一个启发性的陈述能把这两个定理包含在内, 其中前一个暂时还只是猜想. 这里似乎缺少一个新思想, 正如在代数几何中, 要超越射影性假设证明 Riemann--Roch 定理, 也还缺少一个新思想一样 (参见\hyperlink{sga6.e14.sec.2}{第十四讲第 2 节}).

本讨论班不使用 J.-P. Serre 的 \emph{Algèbre locale. Multiplicités} (Lecture Notes in Mathematics 第 11 册, Springer) 中阐述的局部相交理论; 这里只指出, Serre 的课程显然影响了这一理论在 1957 年的形成. 我们也不使用 1958 年 Chevalley 讨论班第二、三讲 (École Normale Supérieure) 所发展的 Chow 环理论. 那一理论本质上依赖非奇异性假设, 而本讨论班的目的恰在于为一般概形, 甚至一般局部赋环拓扑斯, 建立相交理论. 关于这一点, 参见\hyperlink{sga6.e14.expose}{第十四讲}第 4 节和第 8 节的说明, 那里给出了我们的理论与 Chow 环理论的关系, 并提出了推广后者的问题. 因此可以说, 本讨论班给出了一种连贯而自足的相交理论, 但远没有穷尽相交理论中各种有用乃至不可缺少的观点; 所以还应以这里所引 Serre 和 Chevalley 的讲稿作为补充.

我们把 Grothendieck 1957 年的报告 ``Classes of sheaves and the Riemann--Roch theorem'' 作为\hyperlink{sga6.e00.expose}{第零讲}的附录收入本讨论班; 同年在 Princeton 举行的 Borel--Serre 讨论班及其前引论文都以该报告为基础. 报告概述了非奇异拟射影簇的 Riemann--Roch 定理的两种证明. 第一种目前只在特征零时有效, 但对于浸入给出的结论稍强; 除 Borel--Serre 的工作及本讨论班外, 这种证明尚未见于文献. 这篇报告当然不以本讨论班的其他各讲为前提; 它与\hyperlink{sga6.e00.expose}{第零讲}一样, 还可以作为学习本讨论班的入门, 尤其适合尚不熟悉 Borel--Serre 的读者. 刚才提到的证明几乎不加改变就适用于复解析空间的射影态射; 一旦导出范畴中的“外幂”运算问题得到解决, 它无疑也将适用于任意特征 (参见\hyperlink{sga6.e14.expose}{第十四讲}第 1、2、3 节). 从我们采取的观点来看, 没有研究这一问题, 无疑是本讨论班最严重的缺憾.

读者会发现, 封面上列出的参加者中有两位未出现在本讨论班的目录中. 其中 J.-P. Jouanolou 积极参与了讨论班第一部分的技术准备, 但未能参加口头报告; 特别是, 关于射影丛上局部自由模的类环 $K^\circ$ 结构的重要命题 \hyperlink{sga6.e06.thm.1-1}{VI 1.1} 中的线性无关性断言出自他, 而此前该断言只在底概形具有丰沛可逆模的假设下得到证明. 另一位 J.-P. Serre 作了两次口头报告, 在逻辑上独立于讨论班的其余部分, 因而他愿将其另行发表为一篇论文 (参见 J.-P. Serre, ``Grothendieck groups of split reductive group schemes'', 将刊于 \emph{Publications Mathématiques} 第 34 期).

与 \emph{Séminaire de Géométrie Algébrique du Bois Marie} 的其他部分一样, 缩写 SGA 1 至 SGA 6 指该讨论班的各部分; 缩写后的罗马数字表示讲次, 随后的阿拉伯数字则对应所引那一讲内部的编号. 在本讨论班内部引用时省去 SGA 6, 使用 I 4.9 这样的形式, 意为第一讲命题 4.9. 缩写 EGA 指 J. Dieudonné 与 A. Grothendieck 的 \emph{Éléments de Géométrie Algébrique}.

\sgaStableTarget{sga6.front.table-of-contents}\section*{原书目录}
\begin{longtable}{@{}p{0.10\linewidth}p{0.77\linewidth}r@{}}
\toprule
讲次 & 题名及作者 & 原页 \\
\midrule
0 & 一般概形相交理论的计划纲要, A. Grothendieck & 1 \\
附录 & 层的类与 Riemann--Roch 定理, A. Grothendieck & 20 \\
I & 导出范畴中有限性条件的一般理论, L. Illusie & 78 \\
II & 整体预解的存在性, L. Illusie & 160 \\
III & 相对有限性条件, L. Illusie & 222 \\
IV & 赋环拓扑斯的 Grothendieck 群, L. Illusie & 274 \\
V & $\lambda$-环的一般理论, P. Berthelot & 297 \\
VI & 射影丛的 $K^\circ$: 计算及其推论, P. Berthelot & 365 \\
VII & 正则浸入与爆破的 $K^\circ$ 的计算, P. Berthelot & 416 \\
VIII & Riemann--Roch 定理, P. Berthelot & 466 \\
IX & $K_\bullet$ 群的若干计算, P. Berthelot & 498 \\
X & 固有代数概形的相交形式体系, O. Jussila; 附 A. Grothendieck 的“相交理论中的特化” & 519 \\
XI & 未写成讲稿 & -- \\
XII & Picard 函子的一个相对可表性定理, M. Raynaud, S. Kleiman 执笔 & 595 \\
XIII & Picard 函子的有限性定理, S. Kleiman & 616 \\
XIV & 相交理论中的未解决问题, A. Grothendieck & 667 \\
 & 术语索引 & 691 \\
 & 记号索引 & 696 \\
\bottomrule
\end{longtable}
\clearpage
